\begin{pastexam}[label={px:ust-115-ee-ds-06}]{臺聯大電機類 115 學年度 第 6 題}
（單選；2 分）

Given an input sequence of 1, 2, 3 (pushed into the stack in that order), which of the following is
``not possible'' the order in which the elements are popped out?

(A)~3, 2, 1 \qquad (B)~1, 2, 3 \qquad (C)~2, 3, 1 \qquad (D)~3, 1, 2
\end{pastexam}

\begin{pastsoln}
答案是 (D)。

逐項用 push、pop 模擬（資結書第~6.1.5~節）：
\begin{itemize}
\item (A) 3, 2, 1：push 1、push 2、push 3、pop 3、pop 2、pop 1。 做得到。
\item (B) 1, 2, 3：push 1、pop 1、push 2、pop 2、push 3、pop 3。 做得到。
\item (C) 2, 3, 1：push 1、push 2、pop 2、push 3、pop 3、pop 1。 做得到。
\item (D) 3, 1, 2：要先 pop 3，就得先 push 1、2、3，pop 3 之後堆疊由下而上是 $1, 2$，頂端是
$2$。 下一個能 pop 的只有 $2$，不可能先輸出 $1$。 做不到。
\end{itemize}
用資結書定理~6.1 也能直接判斷：輸出順序做不到，若且唯若裡面有「大、小、中」的樣子
（$i < j < k$ 且 $p_j < p_k < p_i$）。 $3, 1, 2$ 本身就是這個樣子，其他三個都沒有。
三個數共有 $3! = 6$ 種排列，做得到的有 $C_3 = 5$ 種，唯一做不到的就是 $3, 1, 2$。
\end{pastsoln}
